\documentclass[10pt,a4paper]{article}
\usepackage[UTF8]{ctex}
\usepackage{amsmath,amssymb,geometry,xcolor,listings,tcolorbox,multicol,enumitem}
\geometry{margin=1.2cm,top=1cm,bottom=0.8cm}
\setlength{\parskip}{0pt}\setlength{\parsep}{0pt}\setlength{\itemsep}{0pt}
\setlist{nosep,leftmargin=*}

\lstset{language=C,basicstyle=\ttfamily\scriptsize,keywordstyle=\color{blue}\bfseries,
  numbers=none,frame=single,breaklines=true,showspaces=false,tabsize=2,
  aboveskip=2pt,belowskip=2pt}

\newtcolorbox{keybox}[1][]{colback=yellow!3!white,colframe=yellow!70!black,
  title=\textbf{考点},fonttitle=\bfseries\small,size=small,top=2pt,bottom=2pt,#1}
\newtcolorbox{bugbox}[1][]{colback=red!3!white,colframe=red!60!black,
  title=\textbf{易错点},fonttitle=\bfseries\small,size=small,top=2pt,bottom=2pt,#1}

\title{\textbf{哈希表 (Hashing) — 考前速查}}
\author{}\date{}

\begin{document}
\maketitle\vspace{-8pt}

\begin{keybox}
\textbf{核心突破：} 不通过比较查找，$O(1)$ 定位。\\
基于比较的查找下界 $\Omega(\log N)$，哈希突破了它。
\end{keybox}

\begin{multicols}{2}

% ====== 概念 ======
\section*{1. 基本概念}

\begin{tabular}{|l|l|}
\hline
$b$ (buckets) & 桶数量 \\
$s$ (slots/bucket) & 每桶槽位数 \\
$n$ & 当前元素数 \\
$T$ & 所有可能 key 总数 \\
$\lambda=\frac{n}{s\cdot b}$ & 负载密度 \\
\hline
\end{tabular}

\vspace{4pt}
\textbf{冲突 (Collision)：} $f(i_1)=f(i_2)$，两不同 key 映射到同桶。\\
\textbf{溢出 (Overflow)：} 桶已满又来了新 key。\\
当 $s=1$ 时，冲突 = 溢出。

\textbf{例子：} 10个C函数，b=26,s=2,$\lambda=10/52=0.19$。\\
hash(x)=首字母-'a'。acos/atan 都在桶0。平均查找=3.73。

% ====== 哈希函数 ======
\section*{2. 哈希函数}

\textbf{要求：} 计算快 + 分布均匀（无偏）。

\textbf{整数：}
\begin{lstlisting}
h(x) = x % TableSize
\end{lstlisting}
\textbf{TableSize 必须质数！} 否则后几位相同者扎堆。

\textbf{字符串 — Hash3（推荐）：}
\begin{lstlisting}
int hash(const char* key, int T) {
    unsigned int h = 0;
    while (*key) h=(h<<5)+*key++;  // h*32+下一字符
    return h % T;
}
\end{lstlisting}
用 32 因为 $32=2^5$，移位比 *27 快。\\
\textbf{长字符串：} 可精选部分字符，早期字符被左移出去丢失影响。

% ====== 分离链接 ======
\section*{3. 分离链接法}

每桶挂一个链表。冲突的串在一起。

\begin{lstlisting}
typedef struct Node { int key; struct Node* next; } Node;
Node* table[T];

void insert(int key) {
    int h = key % T;
    Node* node = malloc(sizeof(Node));
    node->key = key;
    node->next = table[h];  // 头插法 O(1)
    table[h] = node;
}

Node* find(int key) {
    Node* p = table[key % T];
    while (p && p->key != key) p = p->next;
    return p;
}
\end{lstlisting}

\textbf{缺点：} 需要额外指针空间，缓存不友好（内存不连续）。

\columnbreak

% ====== 开放地址 ======
\section*{4. 开放地址法}

冲突后按探测序列找下一个空位：\\
$h_i = (hash(key) + f(i)) \bmod \text{TableSize}$

\subsection*{线性探测 $f(i)=i$}

每次+1。\textbf{一次聚集：} hash 相同的人扎堆。\\
\textbf{期望探测次数：} 插入失败 $\frac12(1+1/(1-\lambda)^2)$

\subsection*{二次探测 $f(i)=i^2$}

+1,+4,+9,+16...\\
\textbf{定理：} 若 TableSize 是质数且表 $>1/2$ 空，一定找到空位。\\
（前 T/2 个 $i^2 \bmod$ 质数互不相同）

\textbf{二次聚集：} 同 hash 者探测路径一样。

\textbf{代码：}
\begin{lstlisting}
int find(int T[], int key, int size) {
    int h = key % size;
    for (int i=0; i<=size/2; i++) {
        int p = (h + i*i) % size;
        if (T[p]==EMPTY) return -1;
        if (T[p]==key) return p;
    }
    return -1;
}
\end{lstlisting}

\textbf{插入优化公式：} $f(i)=f(i-1)+2i-1$（避免每次算 $i^2$）。\\
\textbf{懒惰删除：} 加 flag 标记删除，不真正清空。

\subsection*{双哈希 $f(i)=i\cdot h_2(key)$}

\textbf{推荐：} $h_2(x)=R-(x\%R)$，R 为小于 TableSize 的质数(如 7)。\\
步长取决于第二个哈希函数，基本消除聚集。

\subsection*{三种探测对比}

\begin{tabular}{|c|c|c|}
\hline
& \textbf{聚集} & \textbf{性能} \\
\hline
线性 & 一次聚集严重 & 最差 \\
二次 & 二次聚集 & 中等 \\
双哈希 & 基本无 & 最好 \\
\hline
\end{tabular}

% ====== Rehashing ======
\section*{5. 再哈希 (Rehashing)}

\textbf{何时触发：}
\begin{itemize}
\item 插入失败（没空位）
\item $\lambda$ 超阈值（如 $>0.5$）
\item 查找性能严重下降
\end{itemize}

\textbf{步骤：}
\begin{enumerate}
\item 建更大的表（约 2 倍，质数）
\item 旧表全部元素重哈希插入新表
\item 释放旧表
\end{enumerate}

$O(N)$ 一次，但\textbf{摊还到每次插入 $O(1)$}。\\
\textbf{交互系统注意：} 不能一次性停太久——可逐步迁移。

% ====== 分离链接 vs 开放地址 ======
\section*{6. 分离链接 vs 开放地址}

\begin{tabular}{|c|c|c|}
\hline
& \textbf{分离链接} & \textbf{开放地址} \\
\hline
实现 & 链表 & 纯数组 \\
空间 & 额外指针 & 无额外开销 \\
满后 & 无上限 & Rehashing \\
$\lambda$ & 可 >1 & 必须 <1 \\
缓存 & 差 & 好（连续） \\
\hline
\end{tabular}

% ====== 线性探测找入度 ======
\section*{7. 线性探测反推输入顺序}

\begin{keybox}
\textbf{给定哈希表状态，反推插入顺序。} 线性探测下，元素从 hash 到实际位置之间经过的所有人都比它先插入。

\textbf{入度公式：} $indeg[i] = (i - hash(a[i]) + N) \% N$

\textbf{然后拓扑排序：} 每次选入度为 0 且值最小的输出。
\end{keybox}

\textbf{例子：} N=5, 表=[10,6,15,0,8]。\\
a[4]=8,h=3,in=1(依赖3); a[3]=0,h=0,in=3(依赖0,1,2)...\\
最小候选优先输出。

\end{multicols}
\end{document}
